\chapter*{Resumen}
El mercado de los juegos de cartas coleccionables (TCG) ha experimentado un crecimiento masivo, pero sigue enfrentándose a desafíos críticos como la proliferación de falsificaciones y la falta de confianza en las transacciones entre particulares (\textit{Peer-to-Peer}). Este trabajo presenta \textbf{TraDeck}, una aplicación descentralizada (dApp) diseñada para garantizar la autenticidad y la seguridad en el intercambio de activos digitales.

La solución se apoya en la tecnología \textit{Blockchain} de Ethereum, utilizando el estándar \textbf{ERC-721} para la tokenización de cartas únicas y el estándar \textbf{ERC-20} para la gestión de una economía interna estable. Para optimizar los costes de almacenamiento, se ha implementado una arquitectura de persistencia híbrida que delega la custodia de metadatos e imágenes a la red descentralizada \textbf{IPFS} \textit{InterPlanetary File System}. El núcleo del proyecto reside en sus contratos inteligentes, que integran un protocolo de custodia (\textit{Escrow}) para compraventas seguras e intercambios atómicos (\textit{Atomic Swaps}) que eliminan el riesgo de fraude. El resultado es una plataforma funcional y accesible que demuestra la viabilidad de la cadena de bloques como una base de datos inmutable para el coleccionismo moderno.

\vspace{0.4cm}

\noindent \textbf{Repositorio oficial:} \url{https://github.com/rafseggom/TraDeck} \\
\noindent \textbf{Vídeo demostrativo:} \url{https://youtu.be/GzI2yNninhk}

\vspace{0.6cm}

\textbf{Palabras clave:} Blockchain, Ethereum, Smart Contracts, NFT, IPFS, TCG, Base de Datos Distribuida, Web3.